\documentclass[11pt]{article}
\usepackage{amsmath,amssymb}
\usepackage[margin=1in]{geometry}
\usepackage{hyperref}
\emergencystretch=3em

\title{The Taylor tree for absolutely summable coefficient families (Taylor-tree Stage 4)}
\author{Claude, with Daniel Murfet}
\date{2026-09-07}

\begin{document}
\maketitle
\sloppy

\begin{abstract}
Units 241--247 complete Stage~4 of the Taylor-tree programme (Astra \#28): the polynomial theorem of
Stage~3 is extended to phase and amplitude given on the unit box by \emph{absolutely summable power
series} $\xi(u)=\sum_\gamma c^\xi_\gamma u^\gamma$, $\eta(u)=\sum_\gamma c^\eta_\gamma u^\gamma$ with
$\sum_\gamma|c_\gamma|<\infty$ --- the hypothesis implied by the paper's holomorphy on a polydisc
$D_R$, $R>b=1$, through Cauchy estimates. For $\beta>0$, $k_i>0$, every cutoff $L>0$ and every $N\ge1$,
\[
\Bigl|Z(N)-\sum_{\mu\in\Lambda_L}N^{-\mu}\sum_{j\le d-1}A_{\mu,j}(\log N)^j\Bigr|
\le C\bigl(\xi(0),\|c^\eta\|_1,\|c^\xi\|_1\bigr)\,N^{-L}(1+\log N)^{d-1}
\]
(Headline~XXVII, \texttt{familyTaylorTree\_cutoff\_bound}), with cutoff-free coefficients $A_{\mu,j}$
vanishing off $\Lambda(h,k)$ and a constant depending on the data only through $\xi(0)$ and the two
coefficient masses; hence $O(N^{-L}(\log N)^{d-1})$ and $o(N^{-L'})$ for $L'<L$. The route is
\emph{completion by density}, not a redo of Stage~3 for infinite families: box truncations are
genuine monomial lists, the truncated coefficients form a Cauchy sequence by a Lipschitz
\emph{stability gate}, the truncated integrals converge by dominated convergence at fixed $N$, and
the uniform polynomial bound passes to the limit. Numerics ($d=1$, $\xi=0.3e^{u/2}$,
$\eta=(1-u/3)^{-1}$, $L=5/2$): error$\times N^{5/2}\approx0.05$--$0.06$ for $N=10,\dots,640$. No
\texttt{sorry} and no additional \texttt{axiom} declarations; 9123 jobs.
\end{abstract}

\section{Architecture}

\paragraph{List algebra (u241, \texttt{MonoRepPerm}).} The coefficient functionals of Stage~3 depend
on a monomial list only through its multiset of $(\gamma,c)$ pairs: invariant under permutation,
additive under \texttt{++}; the list product distributes over \texttt{++} (exactly on the left, up to
permutation on the right). The \textbf{power-difference estimate}
$\mathrm{pow}(J\mathbin{+\!\!+}\Delta)^p\sim\mathrm{pow}(J)^p\mathbin{+\!\!+}R_p$ with
$\|R_p\|_1\le p\|\Delta\|_1(\|J\|_1+\|\Delta\|_1)^{p-1}$ is the list form of
$\mathrm{mass}(x^p-y^p)\le pB^{p-1}\mathrm{mass}(x-y)$.

\paragraph{Stability gate (u242, \texttt{CoeffStability}).} If $\eta'\sim\eta\mathbin{+\!\!+}\Delta\eta$,
$J'\sim J\mathbin{+\!\!+}\Delta$, $\xi'(0)=\xi(0)=a$, $\|\eta\|_1\le E$, $\|J\|_1+\|\Delta\|_1\le B$, then for
$\mu>0$
\[
|A_{\mu,j}(\xi',\eta')-A_{\mu,j}(\xi,\eta)|
\le K_kD\bigl(E\beta M_{\mu+1/2,n}(a+B)\,\|\Delta\|_1+M_{\mu,n}(a+B)\,\|\Delta\eta\|_1\bigr),
\]
where the factor $p$ from the power difference is absorbed by shifting the phase order,
$M_{\nu,n,p+1}=M_{\nu+1/2,n,p}$, and the phase orders are summed by the Tonelli identity of unit 234.
Astra \#28 named this the single riskiest step of Stage~4; it passed in one unit.

\paragraph{Coefficient families (u243, \texttt{CoeffFamily}).} Families $c:\mathbb N^d\to\mathbb R$ with
$\sum|c_\gamma|<\infty$; mass, evaluation ($|\mathrm{eval}\,c\,u|\le\mathrm{mass}\,c$ on the cube), box
truncations $\{\gamma:\gamma_i\le m\}$ as monomial lists keeping the constant term; the tail mass
$\sum_{\gamma\notin\mathrm{box}_m}|c_\gamma|\to0$ (Tannery); successive truncations differ by an
appended list of mass at most the tail mass; uniform convergence on the cube.

\paragraph{Integral convergence (u244).} $Z_m(N)\to Z(N)$ at fixed $N\ge0$ by dominated convergence
with the constant bound $\mathrm{mass}(\eta)\,e^{\beta\sqrt N\mathrm{mass}(\xi)}$ on the box.

\paragraph{Coefficient convergence (u245).} $|A(m')-A(m)|\le C_1\,\mathrm{tail}_\xi(m)+C_2\,\mathrm{tail}_\eta(m)$
for $m\le m'$, so the truncated coefficients are Cauchy; $A_{\mu,j}(c^\xi,c^\eta)$ is their limit (at
$\mu\le0$ and off $\Lambda(h,k)$ all terms vanish, so the limit does).

\paragraph{Uniform constants (u246).} $M_{\nu,r}(b)$ is increasing in $b$ and the tail constant in
$|b|$, so the polynomial cutoff constant is bounded by an explicit function of $(\xi(0),E,B)$ whenever
$\|\eta\|_1\le E$, $\|J\|_1\le B$.

\paragraph{Assembly (u247).} Fix $N\ge1$; apply the uniform polynomial bound to every truncation and
let $m\to\infty$ on both sides (\texttt{le\_of\_tendsto'}). No exchange of the asymptotic limit with the
truncation limit is needed.

\section{Correspondence with the paper}
Lean's $N$ is the sample size $n$; the dimension is $d=n_{\rm Lean}+1$; $(\log N)^j$ carries the
paper's $m=j+1$. The hypothesis is the coefficient-summability form of the analytic hypothesis of
\texttt{thm:TaylorTree}; the Cauchy-estimate bridge from holomorphy on $D_R$ to $\sum|c_\gamma|<\infty$ is
a separate (unformalised) obligation, as is $b\neq1$ (scaling
$Z_b(N)=b^{|h|+d}Z_1(Nb^{2|k|};\xi(b\cdot),\eta(b\cdot))$) and the derivative dictionary
$\beta^p\!\int t^{\mu-1}(-\log t)^i(\sqrt t)^pe^{-\beta t+\beta a\sqrt t}=(-\partial_\mu)^i\partial_a^pS_\mu(a)$.

\section{Lessons}
Complete by density: keep Stage~3 for lists, extend only the coefficient algebra and functionals.
Avoid the infinite Cauchy product ($\mathrm{Fin}\,d\to\mathbb N$ has no \texttt{HasAntidiagonal}
instance): truncation differences are literal appended lists, so permutation invariance replaces
convolution algebra. Lean: \texttt{tendsto\_tsum\_of\_dominated\_convergence} lives in
\texttt{Mathlib.Analysis.Normed.Group.Tannery} (import it explicitly); use
\texttt{List.perm\_of\_nodup\_nodup\_toFinset\_eq} with \texttt{List.mem\_toFinset} rather than
\texttt{Finset.toList\_toFinset}; \texttt{limUnder} + \texttt{cauchySeq\_tendsto\_of\_complete} for a limit
defined without choosing a representative; \texttt{cauchySeq\_of\_le\_tendsto\_0} wants a bound at the
smaller index, so tail masses must be shown antitone.
\end{document}
